\documentclass[11pt, a4paper]{article}
\usepackage[margin=2.5cm]{geometry}
\usepackage{amsmath, amssymb, amsthm}
\usepackage{bm}
\usepackage{microtype}
\usepackage{hyperref}
\usepackage{xcolor}
\usepackage{booktabs}
\usepackage{tcolorbox}
\usepackage{listings}
\usepackage[T1]{fontenc}
\usepackage{lmodern}

\definecolor{copper}{HTML}{8f3e22}
\definecolor{ivory}{HTML}{f7f3ea}
\hypersetup{colorlinks=true, linkcolor=copper, urlcolor=copper, citecolor=copper}

\theoremstyle{definition}
\newtheorem{definition}{Definition}[section]
\newtheorem{proposition}{Proposition}[section]

\lstset{language=Python, basicstyle=\small\ttfamily, keywordstyle=\color{copper}\bfseries,
  frame=single, backgroundcolor=\color{ivory}, breaklines=true, columns=fullflexible}

\title{\color{copper}\textbf{Diffusion Transformers + Video \& World Models}\\[4pt]
  \large TLH-3 · Thursday Learning Hours · Applied Models Track}
\author{Thursday Learning Hours}
\date{April 24, 2026}

\begin{document}
\maketitle

\begin{tcolorbox}[colback=ivory, colframe=copper, title=\textbf{Core Thesis}]
The U-Net was a pragmatic choice — not a principled one. Vision Transformers process patches
with global attention from the first layer, follow power-law scaling with compute, and extend
naturally to video by treating time as another patch dimension. DiT proved this works.
Sora showed how far it scales.
\end{tcolorbox}

\tableofcontents
\newpage

\section{DiT: Diffusion Transformer}

Peebles \& Xie~\cite{peebles2023dit} replace the U-Net with a Vision Transformer backbone.

\subsection{Patchification}

An image $\bm{x}_t \in \mathbb{R}^{H \times W \times C}$ is partitioned into non-overlapping
patches of size $p \times p$, then projected to a $d$-dimensional hidden space:
\begin{equation}
  \text{tokens} = \text{Linear}(\text{reshape}(\bm{x}_t, \tfrac{H}{p}, \tfrac{W}{p}, p^2 C)) \in \mathbb{R}^{N \times d}
\end{equation}
where $N = (H/p)(W/p)$ is the number of spatial tokens.

\subsection{Adaptive Layer Norm (adaLN)}

Time and class conditioning via learnable per-token scale and shift:
\begin{equation}
  \text{adaLN}(x, t, c) = \gamma(t, c) \odot \text{LayerNorm}(x) + \delta(t, c)
\end{equation}
where $\gamma, \delta \in \mathbb{R}^d$ come from a shared MLP over the concatenated time+class embedding.

\subsection{DiT Block}

\begin{equation}
  x \leftarrow x + \text{gate}_1 \odot \text{Attn}(\text{adaLN}(x, t, c))
\end{equation}
\begin{equation}
  x \leftarrow x + \text{gate}_2 \odot \text{MLP}(\text{adaLN}(x, t, c))
\end{equation}

After $L$ blocks: linear head + unpatchify → predicted noise $\hat{\bm{\epsilon}}$.

\subsection{Scaling Law}

Peebles \& Xie observe FID $\propto \text{GFLOPs}^{-\alpha}$ — a clean power law — consistent
with LLM scaling~\cite{kaplan2020scaling}. U-Net counterparts plateau.

\begin{center}
\begin{tabular}{llll}
\toprule
Model & Params & GFLOPs & FID (ImageNet 256) \\
\midrule
DiT-S/2 & 33M & 6 & 68.4 \\
DiT-B/2 & 130M & 23 & 43.5 \\
DiT-L/2 & 458M & 80 & 23.3 \\
DiT-XL/2 & 675M & 118 & 2.27 \\
\bottomrule
\end{tabular}
\end{center}

\section{Video Diffusion Models}

\subsection{Factorized Space-Time Attention}

Ho et al.~\cite{ho2022video} model video as a 4D tensor $V \in \mathbb{R}^{T \times H \times W \times C}$
with factorized attention:
\begin{enumerate}
  \item Spatial attention: each frame attends over its own $(H, W)$ tokens
  \item Temporal attention: each spatial position attends over its $T$ temporal instances
\end{enumerate}

\subsection{Spacetime Patches (Sora)}

Brooks et al.~\cite{brooks2024sora} extend patchification to spacetime:
\begin{equation}
  \text{video token}: \Delta t \times p \times p \text{ spacetime patch} \to d\text{-dim vector}
\end{equation}

A video of $T$ frames with patch size $(p, p)$ and temporal stride $\Delta t$ gives
$N = \frac{T}{\Delta t} \cdot \frac{H}{p} \cdot \frac{W}{p}$ total tokens. The same
transformer architecture handles arbitrary $(T, H, W)$ without modification.

\section{World Models via Diffusion}

\subsection{Definition}

A world model approximates $p(s_{t+1} | s_t, a_t)$ — the distribution over next states
given the current state and action.

\subsection{DIAMOND}

Alonso et al.~\cite{alonso2024diamond} use DDPM to model this distribution:
\begin{equation}
  p_\theta(s_{t+1} | s_t, a_t) = \text{DDPM reverse}(\text{noise}; s_t, a_t)
\end{equation}

Action conditioning via cross-attention: $s_t$ is context (key/value tokens); noise is query.
This is identical to text-conditional image generation — ``text'' = $(s_t, a_t)$,
``image'' = $s_{t+1}$.

\subsection{Advantages Over Deterministic Models}

\begin{itemize}
  \item Multimodal: samples from $p(s_{t+1}|\cdot)$ — captures stochastic environments
  \item High-fidelity: diffusion generates more realistic observations than RSSM
  \item Composable: can generate from imagination for model-based RL planning
\end{itemize}

\section{Key Takeaways}

\begin{enumerate}
  \item DiT proves ViT scales better than U-Net — FID power law vs plateau.
  \item Spacetime patches unify image and video — same transformer, more tokens.
  \item Sora's emergent physics arises from scale, not explicit 3D supervision.
  \item Diffusion world models are multimodal — unlike deterministic/AR alternatives.
  \item SiT (DiT + flow matching) shows both architecture and path choice matter.
\end{enumerate}

\section*{References}
\begin{thebibliography}{9}
\bibitem{peebles2023dit}
  Peebles, W., \& Xie, S. (2023). Scalable Diffusion Models with Transformers (DiT).
  ICCV 2023. arXiv:2212.09748.

\bibitem{dosovitskiy2021vit}
  Dosovitskiy, A., et al. (2021). An Image is Worth 16×16 Words: Transformers for Image
  Recognition at Scale. ICLR 2021. arXiv:2010.11929.

\bibitem{ho2022video}
  Ho, J., et al. (2022). Video Diffusion Models. NeurIPS 2022. arXiv:2204.03458.

\bibitem{blattmann2023svd}
  Blattmann, A., et al. (2023). Stable Video Diffusion. arXiv:2311.15127.

\bibitem{brooks2024sora}
  Brooks, T., et al. (2024). Video Generation Models as World Simulators. OpenAI.

\bibitem{alonso2024diamond}
  Alonso, E., et al. (2024). DIAMOND: Diffusion for World Model. NeurIPS 2024.
  arXiv:2405.12399.

\bibitem{bruce2024genie}
  Bruce, J., et al. (2024). Genie: Generative Interactive Environments. arXiv:2402.15391.

\bibitem{kaplan2020scaling}
  Kaplan, J., et al. (2020). Scaling Laws for Neural Language Models. arXiv:2001.08361.

\bibitem{ma2024sit}
  Ma, N., et al. (2024). SiT: Scalable Interpolant Transformers. arXiv:2401.08740.
\end{thebibliography}

\end{document}
